% OWNER: framing
% STATUS: review
% SOURCE OF TRUTH for definitions: ../../GLOSSARY.md — keep this list in sync.
\chapter*{Wykaz skrótów i symboli}
\addcontentsline{toc}{chapter}{Wykaz skrótów i symboli}
\ChapterStatus{draft}{framing}

\begin{description}[style=nextline]
  \item[\idaTermHead{AI}] Artificial Intelligence / sztuczna inteligencja
  \item[\idaTermHead{ASP}] Akademia Sztuki w Szczecinie
  \item[\idaTermHead{Big Five}] Model Wielkiej Piątki cech osobowości
  \item[\idaTermHead{ComfyUI}] Środowisko grafów do eksperymentów z modelami generatywnymi (pracownia autora; nie runtime IDA)
  \item[\idaTermHead{ControlNet}] Conditioning struktury obrazu (depth, krawędzie itd.) --- warstwa pracowni
  \item[\idaTermHead{FLUX}] Rodzina modeli dyfuzyjnych używana we wczesnej fazie (Modal) --- genealogia, nie stan obecny produkcji
  \item[\idaTermHead{GCP}] Google Cloud Platform
  \item[\idaTermHead{HCI}] Human--Computer Interaction
  \item[\idaTermHead{IAT}] Implicit Association Test
  \item[\idaTermHead{IDA}] AI Interior Design Dialogue Research Platform
  \item[\idaTermHead{IPIP-NEO-120}] International Personality Item Pool (120 pozycji, 30 facetów)
  \item[\idaTermHead{LLM}] Large Language Model
  \item[\idaTermHead{PRS}] Perceived Restorativeness Scale
  \item[\idaTermHead{RQ}] Research Question / pytanie badawcze (w pracy: RQ1--RQ6)
  \item[\idaTermHead{RtD}] Research through Design / badania przez projektowanie
  \item[\idaSamHead] Segment Anything Model --- segmentacja w pracowni generatywnej
  \item[\idaTermHead{Vertex AI}] Platforma Google Cloud do modeli generatywnych; produkcyjna generacja obrazu IDA (Gemini Image)
  \item[\idaTermHead{VLM}] Vision--Language Model
  \item[\idaTermHead{XAI}] Explainable Artificial Intelligence
\end{description}
